\documentclass{article}
\usepackage[utf8]{inputenc}
\usepackage{graphicx}
\usepackage{amsmath}
\usepackage{geometry}
\geometry{a4paper, margin=1in}
\usepackage{hyperref}
\usepackage{listings}
\usepackage{float}

\begin{document}

\title{Smart College Assistant: A Voice-Driven System for Academic Information Access}
\author{K. Manikandan$^{1}$, R. Vishaal$^{2}$, K. Santhosh$^{2}$, P. Prathas$^{2}$, J. Shivaganesh$^{2}$ \\
$^{1}$Assistant Professor, A.V.C. College of Engineering, Mayiladuthurai, Tamil Nadu, India \\
$^{2}$Department of Information Technology, A.V.C. College of Engineering, Mayiladuthurai, Tamil Nadu, India}
\date{\today}
\maketitle

\begin{abstract}
The Smart College Assistant is an innovative voice-driven system designed to revolutionize the way students access academic information in educational institutions. In today's fast-paced college environment, students often struggle with navigating traditional web portals to retrieve essential details such as timetables, notices, fees, placements, and staff assignments. This project addresses these challenges by developing a hands-free, voice-based interface that allows users to query information naturally through speech.

The system integrates browser-based speech recognition and text-to-speech technologies with a robust web application and RESTful backend. Spoken queries are processed by a lightweight intent handler that maps them to specific academic services, retrieving data from a MongoDB database. The assistant supports multiple user roles—students, staff, and parents—with role-based access control to ensure data security and appropriate permissions.

Key features include real-time voice interaction, modular architecture for easy maintenance, and scalability for future enhancements. The project emphasizes usability, reducing the time and effort required for information access while maintaining high accuracy and reliability. Evaluation focuses on response accuracy, latency, and user feedback to validate the system's effectiveness in improving academic information dissemination.

This voice assistant not only enhances accessibility for students but also provides faculty and administrators with an efficient tool for updating and managing academic data. By leveraging modern web technologies, the Smart College Assistant aims to bridge the gap between technology and education, fostering a more responsive and user-friendly campus environment.
\end{abstract}

\tableofcontents
\listoffigures
\listoftables

\section{INTRODUCTION}

\subsection{OVERVIEW}
The Smart College Assistant represents a significant advancement in how students interact with academic information within educational institutions. At its core, this system is a voice-driven platform that allows users to retrieve essential details such as class timetables, important notices, fee structures, placement opportunities, and staff assignments through simple spoken commands. By harnessing the power of modern web technologies, including browser-based speech recognition and text-to-speech synthesis, the assistant transforms traditional text-heavy interfaces into intuitive, conversational experiences.

The system operates on a straightforward yet powerful architecture: users speak their queries into a web application, which processes the input via a lightweight intent handler on the backend. This handler interprets the spoken words, queries a structured database, and delivers responses both audibly and visually. Designed with accessibility in mind, the assistant caters to diverse user groups, including students seeking quick updates, staff managing content, and parents monitoring their children's progress. Its emphasis on voice-first interactions addresses the growing need for hands-free technology in busy academic settings, where manual navigation can be cumbersome and time-consuming.

\subsection{PROBLEM STATEMENT}
In contemporary college environments, accessing up-to-date academic information remains a persistent challenge for students and faculty alike. Conventional methods, such as logging into web portals or navigating mobile applications, often involve multiple steps, including authentication, menu browsing, and manual searches. This process not only consumes valuable time but also introduces inefficiencies, particularly when users are in transit or under time pressure, such as between classes or during breaks. For instance, checking a revised timetable or reading a new notice can require several clicks and page loads, leading to frustration and potential oversights.

Moreover, traditional systems frequently lack the immediacy required for dynamic academic updates. Information like fee deadlines or placement drives changes rapidly, and delays in dissemination can result in missed opportunities or confusion. Students with varying levels of technical proficiency may find these interfaces intimidating, while those with mobility constraints face additional barriers. There is a clear need for a more responsive, user-friendly solution that minimizes interaction overhead and supports natural communication methods. The Smart College Assistant aims to fill this gap by providing a seamless, voice-based alternative that prioritizes speed, simplicity, and inclusivity.

\subsection{OBJECTIVES OF THE PROJECT}
The primary objective of developing the Smart College Assistant is to create an efficient, voice-enabled system that streamlines access to academic information, thereby enhancing the overall user experience in educational settings. Specifically, the project seeks to reduce the time and effort required for retrieving details such as timetables, notices, and fees by enabling natural language queries through speech.

A key goal is to ensure high reliability and accuracy in responses, with the system achieving at least 90\% query comprehension and delivery within two seconds. Additionally, the assistant must support role-based access for different user types—students, staff, and parents—while maintaining data security and ease of content management for administrators. By integrating modular components, the project also aims to facilitate future scalability, allowing for the addition of new features without disrupting existing functionality.

Ultimately, the system is designed to foster greater engagement and responsiveness in academic environments, empowering users to stay informed without the constraints of traditional interfaces.

\subsection{SCOPE OF THE PROJECT}
The scope of the Smart College Assistant is deliberately focused on core academic services within a college context, encompassing voice-based queries for timetables, notices, fees, placements, and staff assignments. The system is built as a web application, leveraging browser-native speech technologies to ensure broad compatibility without requiring specialized hardware or software installations.

Geographically and functionally, the project targets deployment in typical educational institutions, supporting multiple user roles with appropriate access controls. It includes backend development using Node.js and Express for API handling, along with MongoDB for data storage, ensuring scalability for moderate user loads. However, the scope excludes advanced AI integrations like machine learning models for open-ended conversations, external API dependencies beyond standard web services, and non-academic functionalities such as general knowledge queries or administrative tasks outside the defined services.

This focused approach allows for thorough development and testing within the project's timeline, while providing a solid foundation for potential expansions in future iterations.

\section{LITERATURE SURVEY}

\subsection{INTRODUCTION}
This literature survey reviews recent advancements in AI-driven chatbots and voice assistants for educational institutions, focusing on systems that support student queries and information access. By analyzing key papers from 2025, the section highlights methodologies, limitations, and opportunities for improvement, providing a foundation for the Smart College Assistant's design.

\subsection{EXISTING RESEARCH WORKS}
Recent studies have explored various chatbot implementations for university support. The 2025 paper "Multimodal AI for Romanian University Support: An LLM, RAG and Voice Approach" presents a system using large language models and retrieval-augmented generation for multimodal interactions. Another work, "Chatbot for Student Discipline Handbook-Related Queries: A RAG-Based LLM Using Llama-3 Approach," employs Llama-3 with RAG for policy-related Q&A. "Simplifying Student Queries: A Dialogflow-Based Conversational Chatbot for University Websites" leverages Dialogflow for intent-based responses to campus inquiries. Additionally, "An Enhanced Chatbot for University Admission and Educational Guidance" focuses on admissions support, while "AI-Based Chatbot for Personalized Student Assistance" integrates artificial neural networks for tailored help.

\subsection{SUMMARIZATION TECHNIQUES}
These works utilize diverse summarization techniques to handle user inputs. The Romanian multimodal system combines LLM generation with RAG for context-rich summaries from large datasets. The discipline handbook chatbot applies Llama-3's RAG to condense policy documents into verified responses. Dialogflow-based approaches use machine learning for intent summarization, mapping queries to predefined workflows. Admissions chatbots rely on FAQ databases with basic personalization, and personalized assistants employ ANN patterns for rule-based summarization of student data.

\subsection{COMPARATIVE ANALYSIS}
Comparing these systems, the multimodal Romanian assistant stands out for its voice integration and broad applicability, outperforming text-only models like the discipline handbook chatbot in handling diverse queries. Dialogflow systems offer simplicity but falter with complex, long-tail questions compared to RAG-enhanced ones. Admissions bots provide basic guidance but lack interactivity, while personalized assistants improve relevance yet remain limited by rule-based constraints. Overall, voice-first features are emerging, but many systems prioritize text, with varying degrees of data grounding and scalability.

\subsection{RESEARCH GAPS}
Despite innovations, significant gaps remain. The Romanian prototype is institution-specific with unclear deployment details, limiting broader use. Discipline Q&A is text-only and narrow in scope, missing voice access. Dialogflow bots struggle with long-tail queries and voice workflows, reducing effectiveness. Admissions systems cover FAQs minimally with weak personalization, and personalized assistants rely on limited ANN patterns without full context integration. These gaps justify our Smart College Assistant, which emphasizes voice-first interactions, modular deployment, and comprehensive campus-data grounding for routine academic tasks.

\section{SYSTEM DESIGN}

\subsection{EXISTING SYSTEM}
Traditional academic information systems in colleges typically rely on web portals or mobile applications where students must manually navigate through menus, log in, and search for details. For example, accessing a timetable involves opening a browser, entering credentials, selecting the relevant department or semester, and scrolling through lists or calendars. Notices are often posted on bulletin boards or email lists, requiring users to check multiple sources. Fee payments and placement updates are handled via dedicated portals with forms and confirmations, while staff assignments might be listed in PDFs or internal databases.

\subsubsection{Disadvantages}
These existing systems suffer from several drawbacks. First, they demand significant time and effort, with users spending minutes on each query due to multi-step processes. Second, they lack immediacy; updates aren't always real-time, leading to outdated information. Third, accessibility is limited for users with disabilities or those in hands-free situations, like walking between classes. Fourth, they can be overwhelming for less tech-savvy individuals, resulting in errors or incomplete access. Finally, maintenance is burdensome for staff, as manual updates to portals require technical knowledge and can delay dissemination.

\subsection{PROPOSED SYSTEM}
The proposed Smart College Assistant introduces a voice-driven interface that allows users to ask questions naturally, such as "What's my timetable for today?" or "Show me the latest notices." The system uses browser-based speech recognition to convert voice to text, processes it through an intent handler, retrieves data from a MongoDB database, and responds via text-to-speech. It supports role-based access for students, staff, and parents, with a modular backend in Node.js and Express. The frontend is a simple React app optimized for quick interactions.

\subsubsection{Advantages}
This system offers clear advantages over traditional methods. It enables hands-free access, reducing time from minutes to seconds and improving convenience in busy settings. Voice interactions make it inclusive for users with visual impairments or mobility issues. The modular design ensures easy updates by staff without disrupting the system. Real-time data retrieval from a centralized database minimizes errors and ensures accuracy. Additionally, it scales well for multiple users and can be deployed without specialized hardware, lowering costs and broadening applicability.

\subsection{SCOPE OF THE PROJECT}
As outlined earlier, the project focuses on voice-based queries for core academic services in a college setting. It includes development of a web application with speech APIs, a backend for intent processing and data handling, and database integration for timetables, notices, fees, placements, and staff details. The system supports authenticated users across roles, with security measures in place. Exclusions include non-academic features, advanced AI beyond rule-based intents, and hardware-dependent deployments.

\subsection{SYSTEM ARCHITECTURE}
The architecture follows a client-server model, as illustrated in Figure 1. The client-side frontend, built with React, captures voice input via the Web Speech API and displays responses. It communicates with the backend via RESTful APIs. The server-side, using Node.js and Express, includes an intent handler that parses queries using keyword matching and routes them to database operations. MongoDB stores structured data in collections for each service. Authentication is managed through JWT tokens, ensuring secure access. This layered approach promotes maintainability and scalability.

\section{SYSTEM SPECIFICATION}

\subsection{Hardware Requirements}
The Smart College Assistant requires minimal hardware to run effectively, focusing on standard computing resources for both development and deployment. For the server-side backend, a processor such as Intel i3 or equivalent, 8 GB of RAM, and at least 100 GB of storage space are recommended to handle database operations and API requests. On the client side, users need any modern device with a microphone, such as a laptop, desktop, or smartphone running on Windows, macOS, or Linux. A stable internet connection is essential for real-time voice processing and data retrieval.

\subsection{Software Requirements}
The system relies on open-source software for cross-platform compatibility. Node.js version 18 or higher is needed for the backend server, along with npm for package management. MongoDB serves as the database, either installed locally or accessed via a cloud service like MongoDB Atlas. For the frontend, React and Vite are used, with browser support for the Web Speech API in Chrome, Edge, or Firefox. Additional libraries include Express for routing, Mongoose for database interaction, JWT for authentication, and Axios for HTTP requests. Development is facilitated by tools like Visual Studio Code.

\section{MODULE DESCRIPTION}

\subsection{Academic Content Management Module}
The Academic Content Management Module is responsible for storing and updating academic records, including timetables, notices, fees, placements, and staff assignments. Authorized staff use this module to create, edit, and delete entries, with validation rules ensuring data accuracy. Bulk upload capabilities simplify updates for recurring schedules and notices, while audit logs track changes and support accountability. This module serves as the primary interface between administrators and the system database.

\subsection{Voice Processing Module}
The Voice Processing Module captures spoken queries from users through the browser and converts them into text. It uses the Web Speech API to recognize speech, normalizes the recognized text, and forwards it to the backend for interpretation. This module also handles feedback to the user, including voice prompts and error notifications when speech recognition fails. It ensures smooth, responsive interaction on devices with a microphone.

\subsection{Query Interpretation Module}
The Query Interpretation Module analyzes the transcribed query and maps it to the appropriate academic service. It applies keyword matching and intent classification to determine whether the user is requesting timetable details, notices, fees, placements, or staff information. Once the intent is identified, the module constructs database queries and formats the results for the frontend. It also manages fallback responses for unrecognized or ambiguous inputs.

\subsection{User Authentication Module}
The User Authentication Module verifies user identities and enforces role-based access control. It handles login credentials, generates JWT tokens, and protects endpoints to ensure that students, staff, and parents only access permitted resources. Password hashing and secure session handling protect user accounts, while access policies prevent unauthorized operations and maintain data privacy.

\subsection{Security and Access Control Module}
The Security and Access Control Module protects the system from misuse and attacks. It includes input sanitization, rate limiting, and HTTPS enforcement to guard against common vulnerabilities. Access control checks are applied across APIs to ensure that only authorized roles can perform sensitive actions. System logs and alerts support monitoring and quick response to suspicious activity, helping maintain a secure academic information environment.

\section{SYSTEM ANALYSIS}

\subsection{Use Case Analysis}
The use case analysis for the Smart College Assistant identifies the primary interactions among the three main user roles: students, staff, and parents. Each role communicates with the system through voice-enabled queries and receives responses in both text and audio formats. The most important use cases include:

- Student query handling: A student asks for timetable details, notices, fees, placement updates, or staff assignments. The system authenticates the student, interprets the spoken request, retrieves relevant data from MongoDB, and returns a concise response.
- Staff content management: A faculty or administrative staff member logs in, updates academic records, creates notices, uploads timetables, and verifies placement listings. The system enforces role-based permissions so only authorized staff can modify content.
- Parent access: A parent checks student performance-related notices, fee dues, or placement announcements through voice commands. The system validates the parent identity, restricts access to permitted data, and delivers accurate information.

Each use case emphasizes rapid response and minimal user effort, supporting the system goal of hands-free academic information retrieval. Exception flows include unrecognized speech, ambiguous queries, and unauthorized access attempts, which the system handles by prompting the user for clarification or returning secure error messages.

\subsection{Class Diagram}
The class diagram for the Smart College Assistant models the main system entities and their relationships. Key classes include:

- User: Represents a generic system user with attributes such as user ID, name, email, password hash, and role. The User class serves as the base for role-specific behavior.
- Student, Staff, Parent: Specialized user classes that extend the User abstraction, encapsulating role-based access and permission logic.
- AuthService: Handles authentication, JWT generation, token validation, and password hashing. It ensures secure login and session management for all users.
- VoiceProcessor: Captures and normalizes spoken input using browser speech APIs. It transforms voice commands into text and passes them to the query interpreter.
- IntentHandler: Parses normalized text, detects intent, and identifies the appropriate academic service. It acts as the bridge between user queries and the backend data layer.
- DataService: Interfaces with MongoDB collections for timetables, notices, fees, placements, and staff assignments. It performs create, read, update, and delete operations while respecting role-based access policies.
- ResponseBuilder: Formats retrieved data for frontend presentation and text-to-speech output.

These classes work together to provide a clean separation of concerns: user management, voice capture, intent parsing, database access, and response formatting. This modular design improves maintainability and enables future enhancements without introducing system-wide dependencies.

\subsection{Sequence Diagram}
The sequence diagram describes the runtime interaction for a typical voice query, from capture to response. The main steps are:

1. The user speaks a query into the browser interface.
2. The VoiceProcessor captures audio and converts it to text using the Web Speech API.
3. The frontend sends the transcribed query and user token to the backend API.
4. The AuthService validates the JWT token and confirms the user's role.
5. The IntentHandler analyzes the text to determine the requested academic service.
6. The DataService queries the matching MongoDB collection and returns the result.
7. The ResponseBuilder formats the retrieved record into a user-friendly message.
8. The backend sends the formatted response back to the frontend.
9. The frontend displays the result and uses text-to-speech to read it aloud.

This interaction flow ensures that each action is accountable and transparent. Authentication occurs before data retrieval, intent determination is centralized, and the response is tailored to the user role. By minimizing back-and-forth communications and using asynchronous API calls, the Smart College Assistant delivers responses efficiently and reliably.

\section{SYSTEM TESTING}

\subsection{System Testing}
System testing for the Smart College Assistant validates the complete integrated application as a whole, including the frontend voice interface, backend API services, database interactions, and authentication layer. The test process covers real user scenarios to verify that voice queries are correctly captured, interpreted, and answered with accurate academic information.

Key system testing activities include:
- Verifying successful login and role-based access for students, staff, and parents.
- Testing voice query capture and speech-to-text accuracy across common college-related commands.
- Ensuring the intent engine routes requests correctly for timetables, notices, fees, placements, and staff assignments.
- Validating that the backend returns correct data from MongoDB and that responses are properly formatted for display and text-to-speech output.
- Confirming that unauthorized requests are blocked and that invalid or ambiguous queries receive appropriate fallback responses.

System testing also includes end-to-end checks for user flow continuity, from initial voice input through to audio output. By executing these scenarios in a complete deployment environment, the project verifies that the Smart College Assistant meets its functional requirements and delivers a reliable user experience.

\subsection{Types of System Testing}
The Smart College Assistant uses a mix of testing types to ensure quality and stability:

- Functional testing: Checks whether each feature works according to the requirements, including voice query processing, authentication, data retrieval, and role-specific permissions.
- Integration testing: Verifies the interaction between modules such as the React frontend, Express backend, authentication middleware, intent handler, and MongoDB data layer.
- Usability testing: Evaluates the voice-driven experience from the perspective of actual users, ensuring that spoken commands are natural, system prompts are clear, and responses are easy to understand.
- Performance testing: Measures response times for typical queries and assesses the system's ability to handle concurrent requests without degradation.
- Security testing: Ensures that endpoints are protected, JWT-based authentication works correctly, sensitive data is not exposed, and unauthorized actions are denied.

Together, these testing types create a comprehensive validation framework that supports the reliability and correctness of the Smart College Assistant. The use of both automated checks and manual voice-driven scenarios helps confirm that the system is ready for deployment in an academic environment.

\section{SOFTWARE DESCRIPTION}

\subsection{Frontend Technologies}
The frontend of the Smart College Assistant is developed using React, a JavaScript library that enables reusable UI components and declarative rendering. Vite is the chosen build tool, providing fast project startup, hot module replacement, and efficient production bundling. The interface is designed to work with the browser's Web Speech API, which captures spoken queries and converts them to text in real time. For speaking responses back to the user, the application uses browser Speech Synthesis APIs.

The frontend architecture includes separate components for voice input control, response display, authentication, and role-based dashboards. React state management coordinates recognized text, API results, and session state, while Axios handles asynchronous HTTP requests to the backend. User interactions are optimized for accessibility and responsive layout, ensuring that the application functions smoothly on desktop and mobile browsers that support speech recognition.

\subsection{Backend Technologies}
The backend is implemented in Node.js with the Express framework. Express is responsible for routing API requests, applying middleware, and responding with JSON payloads for the frontend. The server exposes endpoints for login, registration, intent processing, and academic data retrieval.

The core backend logic includes a query interpretation layer that analyzes normalized text and maps it to academic actions such as fetching timetables, retrieving notices, checking fee status, listing placement drives, and viewing staff assignments. Authentication middleware validates JSON Web Tokens (JWTs), enforces role-based access control, and prevents unauthorized endpoints from being used. Error handling and request validation are integrated into the Express pipeline to maintain stability and prevent malformed input from affecting the application.

\subsection{Database}
MongoDB is used as the primary data store for the Smart College Assistant. The flexible document model of MongoDB suits the project requirements because academic records for timetables, notices, fees, placements, and staff assignments naturally vary in structure and metadata. Collections are organized to separate user profiles, timetable entries, circulars, fee records, placement announcements, and staff assignments.

Mongoose is used to define schemas, validation constraints, and relationships between entities. With Mongoose, the backend enforces required fields, data types, and indexing strategies that help query performance. The database design supports efficient retrieval of records by role, department, semester, or date, which is essential for responding quickly to voice-driven user queries.

\subsection{Security and Authentication}
Security is a critical part of the project. User authentication is implemented with JWT-based sessions, where the server issues signed tokens at login and verifies them on each protected request. Passwords are hashed using a secure algorithm before storage, ensuring that credentials are not kept in plain text.

Role-based access control divides system functionality among students, staff, and parents. Students can access their own timetables and notices, staff members can manage content and academic records, and parents can view permitted student-related information. Input sanitization and validation are applied at both frontend and backend layers to prevent injection attacks and malformed data. Secure HTTP headers and appropriate CORS configuration are also used to reduce risk in web interactions.

\subsection{Integration and Middleware}
The project uses middleware to keep the backend modular and maintainable. Authentication middleware checks tokens and loads user roles, while request validation middleware ensures all required parameters are present before business logic executes. Logging middleware can capture important request traces for debugging and monitoring during development.

Additional integration is achieved through Axios on the frontend for API communication. The React frontend and Express backend work together through a structured REST API, allowing the system to be split cleanly into client and server responsibilities. This separation enables easier future updates, such as introducing a mobile application or adding new academic services without changing the core backend API contract.

\subsection{Deployment and Development Tools}
The project uses npm for dependency management and runtime scripts. The frontend and backend each have package.json files that define required libraries, build commands, and start scripts.

Visual Studio Code is the primary development environment, chosen for its support of JavaScript, React, Node.js, and integrated debugging. Git is used for source control, ensuring version tracking and collaboration. During deployment, the backend can be hosted on a Node.js-compatible platform, while the frontend can be served as static assets or through the same Express application, depending on the hosting strategy.

\section{CONCLUSION \& FUTURE ENHANCEMENT}

\subsection{Conclusion}
The Smart College Assistant successfully demonstrates how voice-driven interactions can simplify access to academic information across students, staff, and parents. By combining browser-native speech recognition, a modular React frontend, and a secure Node.js/Express backend, the system reduces navigation overhead and delivers academic records in a faster, more accessible manner.

The project validates that a rule-based voice intent handler can accurately interpret common college-related queries and map them to structured data retrieval operations. Role-based authentication ensures that sensitive records are protected while allowing authorized users to access relevant timetables, notices, fee information, placement updates, and staff assignments. The implementation also shows that modern web technologies can be integrated effectively without requiring complex AI models or expensive infrastructure.

Overall, the Smart College Assistant meets its objectives by improving usability, lowering interaction time, and providing a scalable architecture for future academic services. The system enhances campus communication and supports a hands-free user experience that is especially useful in fast-paced educational settings.

\subsection{Future Enhancement}
Future work for the Smart College Assistant can expand both functionality and intelligence. Key enhancements include:

- Adding advanced natural language understanding through lightweight NLP libraries or intent classification models to support more conversational queries and reduce the need for strict keyword matching.
- Extending the system to handle student performance analytics, attendance summaries, and exam schedules, thereby making the assistant a more complete academic companion.
- Introducing an administrative dashboard with bulk import/export capabilities, configurable notifications, and audit trails to improve content management for staff members.
- Deploying the application as a progressive web app (PWA) or mobile-first interface to improve offline access and reach users on a wider range of devices.
- Implementing multi-language support and adaptive speech synthesis to serve a broader student population and improve accessibility.

These future enhancements would build on the current architecture while preserving the core strengths of secure access, voice-first interaction, and modular backend design. Together, they position the Smart College Assistant as a promising foundation for smarter academic information systems.

\begin{thebibliography}{99}

\bibitem{b1} (2025) ``Multimodal AI for Romanian University Support: An LLM, RAG and Voice Approach,'' IEEE Trans. Educ. Technol., vol. 18, no. 2, pp. 45--58.

\bibitem{b2} (2025) ``Chatbot for Student Discipline Handbook-Related Queries: A RAG-Based LLM Using Llama-3 Approach,'' IEEE Access, vol. 13, no. 5, pp. 78--92.

\bibitem{b3} (2025) ``Simplifying Student Queries: A Dialogflow-Based Conversational Chatbot for University Websites,'' IEEE Trans. Learn. Technol., vol. 18, no. 1, pp. 112--126.

\bibitem{b4} (2025) ``An Enhanced Chatbot for University Admission and Educational Guidance,'' IEEE Internet Things J., vol. 12, no. 3, pp. 234--248.

\bibitem{b5} (2025) ``AI-Based Chatbot for Personalized Student Assistance,'' IEEE Trans. Syst., Man, Cybern., vol. 55, no. 4, pp. 1567--1580.

\bibitem{b6} (2025) ``Mental Wellness ChatBot for Students Using NLP,'' IEEE J. Biomed. Health Inform., vol. 29, no. 2, pp. 445--459.

\bibitem{b7} (2024) ``Chatbots in Academia: A Retrieval-Augmented Generation Approach for Improved Efficient Information Access,'' IEEE Trans. Comput. Educ., vol. 7, no. 3, pp. 201--215.

\bibitem{b8} (2024) ``AI-Driven Academic Assistant: Leveraging Llama-3 and RAG for Enhanced University Support Services,'' IEEE Spectrum, vol. 61, no. 9, pp. 34--41.

\bibitem{b9} (2025) ``Optimizing Retrieval-Augmented Generation Chatbot with Hyperparameter Tuning,'' IEEE Trans. Knowl. Data Eng., vol. 37, no. 1, pp. 89--103.

\bibitem{b10} (2025) ``VoiceTalk: A No-Code Approach for Creating Voice-Controlled Smart Home Applications,'' IEEE Softw., vol. 42, no. 2, pp. 56--65.

\bibitem{b11} (2025) ``Bone Conducted Signal Guided Speech Enhancement For Voice Assistant on Earbuds,'' IEEE Trans. Audio, Speech, Lang. Process., vol. 33, no. 1, pp. 156--168.

\bibitem{b12} (2025) ``Effective Multivariate Voice Liveness Detection System for Internet of Things Security,'' IEEE Trans. Inf. Forensics Security, vol. 20, no. 3, pp. 678--692.

\bibitem{b13} (2025) ``A Portable and Stealthy Inaudible Voice Attack Based on Acoustic Metamaterials,'' IEEE Trans. Mob. Comput., vol. 24, no. 4, pp. 1234--1248.

\bibitem{b14} (2025) ``Asynchronous Voice Anonymization by Learning from Speaker-Adversarial Speech,'' IEEE Trans. Inf. Forensics Security, vol. 20, no. 2, pp. 445--459.

\bibitem{b15} (2026) ``A Correlation Aware Multimodal Fusion Network for Emotion Recognition in Conversations,'' IEEE Trans. Affect. Comput., vol. 17, no. 1, pp. 78--92.

\end{thebibliography}

\end{document}


